% TOP_PROBLEM: {"id": 1224, "title": "Borsuk's Conjecture", "category_id": 6, "category": {"id": 6, "name": "geometry", "display_name": "Geometry", "description": "Euclidean and non-Euclidean geometry, geometric structures.", "slug": "geometry", "order_index": 6, "created_at": "2026-07-31T15:26:25.671Z"}, "status": "open"}
% FIRST_POSTED: null
% ENTRY_KIND: statement_only
% !TeX program = lualatex
\documentclass[11pt]{article}
\usepackage[margin=25mm]{geometry}
\usepackage{fontspec}
\setmainfont{Latin Modern Roman}
\usepackage{amsmath,amssymb,mathtools}
\usepackage{unicode-math}
\setmathfont{Latin Modern Math}
\usepackage{xurl,hyperref}
\hypersetup{colorlinks=true,urlcolor=blue}
\setcounter{secnumdepth}{0}
\setlength{\parindent}{0pt}
\setlength{\parskip}{5pt}
\setlength{\emergencystretch}{3em}
\newcommand{\R}{\mathbb R}
\newcommand{\N}{\mathbb N}
\newcommand{\Z}{\mathbb Z}
\newcommand{\Q}{\mathbb Q}
\newcommand{\C}{\mathbb C}
\newcommand{\D}{\mathbb D}
\newcommand{\F}{\mathbb F}
\newcommand{\sref}[2]{\hyperlink{source:#1:#2}{[S#2]}}
\newcommand{\cataloguescope}[1]{#1}
\begin{document}
% UnsolvedMath ID 1224; source catalogue code GEO-029
\setcounter{section}{1223}
\section{Borsuk's Conjecture}\label{1224}
{\small UnsolvedMath ID 1224 · Geometry}
\subsection{Definitions and mathematical statement}

\paragraph{Shared notation.}$\mathbb N$, $\mathbb Z$, $\mathbb Q$, $\mathbb R$, and $\mathbb C$ denote the natural numbers, integers, rationals, reals, and complex numbers, respectively. All additional parameters and hypotheses are specified in the question.


\paragraph{Statement recorded in the supplied source.}
For every $n\geq1$, can each bounded subset of $\mathbb R^n$ with positive diameter be partitioned into at most $n+1$ subsets whose diameters are all strictly smaller than that of the original set?
\par\smallskip\textit{Formulation references:} \sref{1224}{1}, \sref{1224}{2}.
\subsection{Short English statement}
For every $n\geq1$, can each bounded subset of $\mathbb R^n$ with positive diameter be partitioned into at most $n+1$ subsets whose diameters are all strictly smaller than that of the original set?
\subsection{Sources}
\begin{itemize}
\item[\hypertarget{source:1224:1}{S1}] ULAM.AI and the UnsolvedMath Contributors, UnsolvedMath: A Curated Collection of Open Mathematics Problems, record 1224 (GEO-029). Catalogue attribution under CC BY 4.0.
\url{https://huggingface.co/datasets/ulamai/UnsolvedMath}.
\item[\hypertarget{source:1224:2}{S2}] Original problem formulation and mathematical context.
\url{https://doi.org/10.1007/BF01448035}.
\end{itemize}
\label{end:1224}
\end{document}
